%% Relatório técnico do Sistema Operacional
\documentclass[
    12pt,
    openright,
    a4paper,
    english,
    brazilian
]{abntex2}

% Pacotes fundamentais
\usepackage{lmodern}
\usepackage[T1]{fontenc}
\usepackage[utf8]{inputenc}
\usepackage{indentfirst}
\usepackage{color}
\usepackage{graphicx}
\usepackage{microtype}
\usepackage{float}
\usepackage{booktabs}
\usepackage{array}
\usepackage{listings}
\usepackage{xcolor}

% Citações no padrão utilizado nos relatórios anteriores
\usepackage[brazilian,hyperpageref]{backref}
\usepackage[num]{abntex2cite}
\citebrackets()

\renewcommand{\backrefpagesname}{Citado na(s) página(s):~}
\renewcommand{\backref}{}
\renewcommand*{\backrefalt}[4]{
    \ifcase #1
        Nenhuma citação no texto.
    \or
        Citado na página #2.
    \else
        Citado #1 vezes nas páginas #2.
    \fi
}

% Configuração para eventuais trechos de código
\definecolor{verde}{rgb}{0,0.5,0}
\lstset{
    language=C,
    basicstyle=\ttfamily\scriptsize,
    keywordstyle=\color{blue},
    stringstyle=\color{verde},
    commentstyle=\color{gray},
    extendedchars=true,
    showspaces=false,
    showstringspaces=false,
    numbers=left,
    numberstyle=\tiny,
    breaklines=true,
    backgroundcolor=\color{green!10},
    breakautoindent=true,
    captionpos=b,
    xleftmargin=0pt
}

% Informações para capa e folha de rosto
\titulo{Projeto e Desenvolvimento de um Sistema Operacional para o Processador MIPS-Lite}
\autor{Guilherme Ferreira Lourenço \\ 143172}
\local{São José dos Campos -- Brasil}
\data{2026}
\instituicao{
    Docente: Prof. Dr. Tiago de Oliveira
    \par
    Universidade Federal de São Paulo -- UNIFESP
    \par
    Instituto de Ciência e Tecnologia -- Campus São José dos Campos
}
\tipotrabalho{Relatório técnico}
\preambulo{Relatório apresentado à Universidade Federal de São Paulo como parte dos requisitos para aprovação na disciplina de Laboratório de Sistemas Computacionais: Sistemas Operacionais.}

\definecolor{blue}{RGB}{41,5,195}

\makeatletter
\hypersetup{
    pdftitle={\@title},
    pdfauthor={\@author},
    pdfsubject={\imprimirpreambulo},
    pdfcreator={LaTeX with abnTeX2},
    pdfkeywords={sistema operacional, processos, memória, arquivos, entrada e saída, MIPS-Lite},
    colorlinks=true,
    linkcolor=blue,
    citecolor=blue,
    filecolor=magenta,
    urlcolor=blue,
    bookmarksdepth=4
}
\makeatother

\setlength{\parindent}{1.3cm}
\setlength{\parskip}{0.2cm}
\makeindex

\begin{document}

\selectlanguage{brazilian}
\frenchspacing

% ELEMENTOS PRÉ-TEXTUAIS
\imprimircapa
\imprimirfolhaderosto*

\begin{resumo}
Este relatório apresenta o desenvolvimento de um projeto de um sistema operacional para um processador MIPS-Lite, contendo gerenciamento de processos, memória, sistema de arquivos e dispositivos de entrada e saída. O sistema consiste em um kernel monolítico escrito em C-, organizado em um único arquivo. O desenho prevê escalonamento \emph{round-robin}, quantum fixo e \emph{swapping} entre as memórias do processador e um disco rígido simulado em Verilog. A memória de instruções e a RAM possuem 4096 palavras. O supervisor previsto para a região protegida ainda não está na imagem. O projeto utiliza até oito diretórios planos e dezesseis executáveis, inicia diretamente no kernel após o \emph{reset} e especifica chamadas de sistema bloqueantes para mediar a E/S dos programas. A opção 1 já executa um programa sem esses mecanismos pendentes.

\noindent
\textbf{Palavras-chave}: sistema operacional. processos. memória. sistema de arquivos. entrada e saída. MIPS-Lite.
\end{resumo}

\pdfbookmark[0]{\listtablename}{lot}
\listoftables*
\clearpage

\pdfbookmark[0]{\contentsname}{toc}
\tableofcontents*

% ELEMENTOS TEXTUAIS
\textual

\chapter{Introdução}

Um sistema computacional é formado por recursos de processamento, memória e dispositivos que precisam ser compartilhados de maneira controlada. Quando cada programa acessa esses recursos diretamente, o controle da execução torna-se dependente da aplicação e aumenta a possibilidade de conflitos. O sistema operacional atua como uma camada intermediária entre os programas e o \emph{hardware}, sendo responsável por organizar o uso da CPU, controlar as áreas de memória, manter informações persistentes e fornecer meios padronizados para comunicação com os periféricos.

O presente trabalho define um sistema operacional para o processador MIPS-Lite implementado em Verilog nas etapas anteriores do projeto. O processador possui arquitetura baseada em MIPS, instruções de 32 bits e memórias separadas para dados e instruções. O software é produzido pelo compilador C- do projeto, que realiza as etapas de análise do código-fonte e gera o binário utilizado na inicialização da memória de instruções.

O contrato do sistema substitui a ROM simples por uma memória de instruções gravável e acrescenta disco rígido simulado, temporizador de \emph{quantum}, mecanismo de interrupção e registradores de controle acessíveis pelo kernel. O sistema operacional ocupa a memória de instruções após o \emph{reset}, enquanto os programas permanecem no disco e são carregados sob demanda.

Além de integrar os componentes desenvolvidos nos laboratórios anteriores, a construção desse sistema permite observar, em uma plataforma reduzida, problemas que normalmente permanecem ocultos por sistemas operacionais de propósito geral. A preservação de contexto, a arbitragem da CPU, a movimentação de imagens entre memória e disco e a mediação da entrada e saída tornam explícita a relação entre arquitetura de computadores, compiladores e software de sistema. O projeto possui, portanto, motivação didática e técnica: validar esses mecanismos em um processador conhecido e compreender os compromissos impostos por recursos limitados.

O relatório apresenta primeiro as técnicas de gerenciamento empregadas por sistemas operacionais e, em seguida, registra a organização e os algoritmos escolhidos para este sistema.

\chapter[Objetivos]{Objetivos}

\section{Geral}

Documentar o projeto de um sistema operacional em C- que gerencia os recursos do processador MIPS-Lite, executa programas armazenados em um disco rígido simulado e oferece uma interface de seleção pelos dispositivos da placa FPGA.

\section{Específicos}

Os objetivos específicos deste relatório são:

\begin{itemize}
    \item estudar técnicas de gerenciamento de processos, memória, sistemas de arquivos e dispositivos de entrada e saída;
    \item definir a organização monolítica do kernel e seus contratos com o processador;
    \item especificar a execução individual e a multitarefa preemptiva com escalonamento \emph{round-robin};
    \item definir o contexto preservado durante a interrupção;
    \item estabelecer o uso do disco simulado para programas e imagens de memória;
    \item definir os diretórios e as operações de listagem, carregamento, renomeação, movimentação e remoção;
    \item registrar o mapa de memória e a organização da entrada e saída.
\end{itemize}

\chapter{Fundamentação Teórica}

\section{Sistemas Operacionais}

O sistema operacional é o componente responsável por controlar a execução dos programas e administrar os recursos disponíveis no sistema computacional. Ele pode ser compreendido como uma máquina estendida, pois oferece abstrações mais simples do que a interface direta do \emph{hardware}, e também como um gerenciador de recursos, pois decide quando e como cada programa poderá utilizar a CPU, a memória e os dispositivos \cite{tanenbaum2016}.

Entre as arquiteturas utilizadas na organização de sistemas operacionais destacam-se os kernels monolíticos, os microkernels e os sistemas em camadas. Em um kernel monolítico, os principais serviços executam em um mesmo espaço e podem chamar uns aos outros diretamente. Essa abordagem reduz a complexidade da comunicação entre módulos, embora exija cuidado para que falhas em um serviço não comprometam o restante do sistema. Em um microkernel, apenas os mecanismos fundamentais permanecem no núcleo, enquanto outros serviços são executados separadamente e se comunicam por troca de mensagens \cite{silberschatz2015}.

Devido às limitações da linguagem C- e ao objetivo didático do projeto, o sistema utiliza um kernel monolítico. Suas funções são organizadas logicamente em seções de processos, memória, arquivos, entrada e saída e interface com o usuário.

\section{Gerenciamento de Processos}

Um processo é um programa em execução, acompanhado do estado necessário para que sua atividade possa ser interrompida e retomada. Além do código do programa, fazem parte desse estado o contador de programa, os registradores, a pilha e os dados utilizados durante a execução \cite{tanenbaum2016}. O sistema operacional mantém essas informações em uma estrutura denominada Bloco de Controle de Processo, ou PCB.

\subsection{Estados de um processo}

No kernel deste projeto, o vetor \texttt{proc\_state} representa os estados dos processos. O valor 0 indica uma posição livre na tabela, o valor 1 identifica um processo pronto e o valor 2 identifica o processo em execução. Ao final de um programa, sua posição é liberada e pode ser utilizada por um novo processo.

\subsection{Escalonamento da CPU}

O escalonador escolhe qual processo pronto deverá utilizar a CPU. Entre os algoritmos conhecidos estão o primeiro a chegar, primeiro a ser atendido, o menor processo primeiro, o escalonamento por prioridades e a alternância circular. Cada técnica favorece uma característica diferente, como simplicidade, tempo médio de espera ou capacidade de resposta \cite{silberschatz2015,stallings2018}.

No escalonamento \emph{round-robin}, os processos prontos são mantidos em uma fila circular. Cada processo executa durante uma fatia de tempo denominada \emph{quantum}. Quando essa fatia termina, o processo é interrompido e inserido novamente no final da fila, desde que ainda não tenha concluído. A técnica é apropriada para sistemas interativos por impedir que um único programa mantenha a CPU indefinidamente \cite{arpaci2018}.

O tamanho do \emph{quantum} influencia o funcionamento do sistema. Um valor muito pequeno aumenta o número de trocas de contexto e o custo de gerenciamento. Um valor muito grande aproxima o comportamento do escalonamento não preemptivo. Na arquitetura adotada, o \emph{quantum} é fixado em 50 instruções concluídas de usuário. Quando existe somente um processo pronto ou executando, o temporizador é desabilitado, pois uma preempção apenas devolveria a CPU ao mesmo processo.

\subsection{Troca de contexto}

A troca de contexto ocorre quando a CPU deixa de executar um processo e passa a executar outro. Para que o primeiro processo possa continuar posteriormente, seu estado deve ser armazenado no PCB ou em uma área reservada pelo sistema. Em seguida, o estado do próximo processo é restaurado.

Os dados normalmente preservados incluem o contador de programa, o ponteiro da pilha, os registradores de propósito geral, o endereço de retorno e registradores especiais. Durante a troca, o kernel também precisa impedir que uma nova interrupção modifique um contexto ainda incompleto. Embora a operação não realize trabalho útil da aplicação, ela é necessária para permitir a multiprogramação.

\section{Gerenciamento de Memória}

O gerenciador de memória controla quais áreas estão ocupadas, associa essas áreas aos processos e impede que a alocação de um programa sobrescreva os dados de outro. Sistemas maiores podem utilizar memória virtual para oferecer a cada processo um espaço de endereçamento próprio, mesmo quando parte dele está armazenada em disco \cite{tanenbaum2016}.

\subsection{Partições fixas e variáveis}

Na divisão por partições fixas, a memória é separada previamente em regiões de tamanho conhecido. Cada processo ocupa uma partição disponível. A implementação é simples, mas pode ocorrer fragmentação interna quando o processo utiliza apenas parte do espaço reservado.

Nas partições variáveis, o tamanho da região é determinado conforme a necessidade do processo. Essa técnica reduz a fragmentação interna, porém cria espaços livres entre áreas ocupadas, conhecidos como fragmentação externa. Algoritmos como \emph{first fit}, \emph{best fit} e \emph{worst fit} podem ser utilizados para escolher um espaço livre.

\subsection{Paginação e segmentação}

A paginação divide o espaço lógico e a memória física em blocos de tamanho fixo. Uma tabela de páginas estabelece a correspondência entre as páginas do processo e os quadros presentes na memória. Como as páginas podem ocupar quadros não contíguos, a fragmentação externa é reduzida. Entretanto, a técnica exige tradução de endereços e estruturas adicionais no \emph{hardware} \cite{arpaci2018,bryant2016}.

Na segmentação, o programa é separado de acordo com unidades lógicas de tamanho variável, como código, dados e pilha. Cada segmento possui uma base e um limite. Essa organização representa melhor a estrutura do programa, mas também pode apresentar fragmentação externa.

\subsection{\emph{Swapping}}

O \emph{swapping} transfere o conteúdo de um processo entre a memória principal e uma unidade de armazenamento secundário. Quando um processo deixa a CPU, sua imagem pode ser gravada no disco; ao voltar a executar, essa imagem é restaurada. A técnica permite manter mais processos conhecidos pelo sistema do que imagens simultaneamente residentes na memória, porém seu custo é maior do que uma simples troca de registradores devido à quantidade de dados transferidos.

\section{Gerenciamento de Sistemas de Arquivos}

Arquivos permitem armazenar informações de forma persistente. O sistema de arquivos define como os dados são identificados, organizados, localizados, lidos e modificados. Para isso, mantém metadados como nome ou identificador, tamanho, localização e estado de ocupação \cite{tanenbaum2016}.

\subsection{Organização de diretórios}

Um diretório de nível único mantém todos os arquivos em uma mesma lista. Sua implementação é simples e adequada quando existem poucos arquivos e apenas um usuário. Também é possível manter vários diretórios planos independentes, sem permitir que um deles contenha outro. Diretórios hierárquicos permitem pastas e subpastas, oferecendo melhor organização, mas exigem tratamento de caminhos e estruturas mais complexas.

\subsection{Métodos de alocação}

Na alocação contígua, os blocos de um arquivo são armazenados em sequência. O acesso sequencial e o acesso direto são simples, pois basta conhecer o bloco inicial e o tamanho. Como desvantagem, o crescimento de um arquivo pode exigir sua movimentação e as remoções podem deixar regiões livres entre os arquivos.

Na alocação encadeada, cada bloco indica a localização do próximo. O arquivo pode crescer utilizando blocos espalhados, porém o acesso aleatório torna-se mais lento. A alocação indexada mantém uma estrutura com os endereços dos blocos do arquivo, facilitando o acesso direto ao custo de espaço adicional para os índices \cite{silberschatz2015,arpaci2018}.

\subsection{Operações sobre arquivos}

As operações fundamentais são criar, ler, atualizar e remover, frequentemente representadas pela sigla CRUD. No sistema definido, a criação corresponde à importação de um executável preparado pelo empacotador; a leitura compreende a listagem e o carregamento; a atualização abrange a troca de seu identificador numérico e sua movimentação entre diretórios; e a remoção desativa a entrada correspondente.

\section{Gerenciamento de Dispositivos de Entrada e Saída}

Os dispositivos de entrada e saída apresentam velocidades e interfaces diferentes. O sistema operacional utiliza controladores e \emph{drivers} para esconder essas particularidades dos programas. Dessa forma, a aplicação solicita uma operação por uma interface conhecida, enquanto o kernel realiza a comunicação com o dispositivo \cite{love2010,stallings2018}.

\subsection{E/S programada e interrupções}

Na E/S programada, a CPU envia um comando ao controlador e consulta repetidamente seu estado até o término da operação. Essa técnica, conhecida como \emph{polling}, é simples, mas mantém o processador ocupado durante a espera.

Com interrupções, o controlador sinaliza a CPU quando a operação termina. O processador pode realizar outro trabalho durante a espera, porém precisa preservar o estado da tarefa interrompida e executar um tratador apropriado.

\subsection{E/S mapeada em memória}

Na E/S mapeada em memória, os registradores dos controladores ocupam endereços do mesmo espaço utilizado pelas instruções de carga e armazenamento. Assim, as operações comuns de leitura e escrita podem iniciar comandos, informar parâmetros e consultar o estado de um dispositivo. O processador deve possuir um decodificador que diferencie acessos à RAM de acessos aos periféricos.

\subsection{Acesso direto à memória}

O Acesso Direto à Memória, ou DMA, permite que um controlador transfira blocos de dados entre um dispositivo e a memória sem exigir que a CPU copie cada palavra. A CPU configura a origem, o destino e o tamanho da transferência. Em seguida, o controlador realiza a cópia e informa sua conclusão por estado ou interrupção \cite{patterson2014}. A técnica é apropriada para transferências entre o disco e as memórias do processador.

\chapter{Definição do Sistema Operacional}

\section{Visão Geral}

O sistema operacional é um kernel monolítico em C- destinado ao processador MIPS-Lite. A linguagem C- utilizada pelo compilador não possui \texttt{\#include}, ligação de múltiplos arquivos, estruturas ou cadeias de caracteres. Por esse motivo, o código do kernel ocupa um único arquivo e suas estruturas são representadas por vetores paralelos.

A interface do menu é o LCD de duas linhas. \texttt{KEY[3]} avança a opção, \texttt{KEY[2]} recua e \texttt{KEY[1]} confirma a tela. Os \emph{switches} e os sete segmentos ficam para a entrada e a saída numérica dos programas, pelas instruções \texttt{IN} e \texttt{OUT}.

O sistema inicia diretamente após o \emph{reset}, com o kernel residente nos endereços 0 a 3071. A folha da vez, 3072 a 3839, recebe um programa. O supervisor previsto para 3840 a 4095 ainda não está na imagem: essa faixa permanece em zero. Quando existir, o supervisor será um mecanismo de transição, não um segundo kernel. Ele não escolhe processos, não altera PCBs, não interpreta o sistema de arquivos, não atende operações de entrada e saída e não controla o menu. Esse papel continua especificado e ainda não está no hardware.

\section{Estado Atual e Alcance da Especificação}

Este capítulo descreve a arquitetura-alvo que deverá orientar a integração entre kernel, compilador e processador. Ele não afirma que todos os componentes já estejam implementados. No estado atual, \texttt{so.cms} ainda descreve quatro PCBs e não chama \texttt{create\_process} na opção 1. Essa opção já copia setores: \texttt{find\_program} e, se falhar, \texttt{show\_error}; \texttt{loadHD(36, 1536)} e \texttt{loadROM(cat\_sector[idx], 3072)}, ambos interrompidos por \texttt{show\_error} quando o retorno não é 0; senão \texttt{lcdshow(3, id, 0)} e \texttt{runat(3072)}, sem \texttt{wait\_choice} antes do programa. \texttt{loadHD}, \texttt{storeHD} e \texttt{loadROM} são emitidos em 4092 a 4095 e devolvem 0 quando \texttt{hd\_status} está concluído e 1 quando é erro. Os endereços 1020 a 1023 não são mais o disco. O processador tem memórias de 4096 palavras, HD simulado e o catálogo do LCD no endereço 4091. A escrita em 4088 arma o retorno só com destino em 3072--3839 e HD descongelado. Temporizador, modo privilegiado, supervisor, \emph{round-robin}, CRUD, leitura do setor 0 e \texttt{IN}/\texttt{OUT} como \emph{trap} continuam ausentes. As opções 2 e 3, o catálogo fixo e as boas-vindas permanecem.

A seção do menu registra o que a simulação já percorre. As seções seguintes empregam o presente normativo para definir o contrato da arquitetura-alvo, e não para registrar resultados experimentais já obtidos. Expressões como ``o kernel realiza'' devem ser entendidas como requisitos de projeto. O futuro é reservado às etapas ainda necessárias para implementar e validar esses requisitos por testes de módulo, integração em simulação e execução na FPGA.

\section{Menu no LCD}
\label{sec:menu-lcd}

O menu exercitado em simulação está em \texttt{so.cms}. O compilador traduz três chamadas para leituras e escritas da palavra no endereço 4091, sem opcode novo na ISA. \texttt{lcdshow(modo, id, pid)} monta e grava a palavra. \texttt{lcdack()} devolve o bit 27, zero ou um. \texttt{lcdsel()} devolve o índice da opção, nos bits 26 a 22. O zero é o primeiro item. O C- não tem operador de deslocamento; as máscaras são emitidas pelo gerador de \emph{assembly}, com \texttt{sl}, \texttt{sr} e \texttt{andi}.

Na escrita, o modo ocupa os bits 7 a 0, o identificador do programa os bits 17 a 8 e o PID os bits 21 a 18. A leitura devolve esses mesmos campos, o índice e o bit de confirmação. Os bits 31 a 28 ficam em zero. Toda escrita apaga a confirmação. Se o modo mudar, o índice volta a zero; se o modo for o mesmo, o índice permanece. \texttt{KEY[1]} liga a confirmação e não mexe no índice, desde que a CPU não esteja parada em \texttt{IN} ou \texttt{OUT}. O pulso que solta uma dessas duas instruções entrega o dado e não confirma o menu. Nesse mesmo ciclo, uma seta não altera o índice. \texttt{KEY[2]} e \texttt{KEY[3]} só andam quando a tela tem mais de um item e a CPU não está nessa espera. \texttt{KEY[0]} é o \emph{reset} da placa, ativo em nível baixo. O menu não o lê.

O painel tem duas linhas de dezesseis caracteres. O texto sai de um catálogo fixo no Verilog. O programa escolhe o modo, o identificador e o PID; não envia caracteres. A \autoref{tab:lcd} lista as telas.

\begin{table}[H]
\centering
\caption{Telas fixas do LCD}
\label{tab:lcd}
\small
\begin{tabular}{clccp{6.2cm}}
\toprule
\textbf{Modo} & \textbf{Nome} & \textbf{Itens} & \textbf{Linha 1} & \textbf{Linha 2} \\
\midrule
1 & MENU & 3 & MENU & \texttt{>1 RUN}, \texttt{>2 MLT} ou \texttt{>3 ARQ} \\
2 & boas-vindas & 1 & BEM-VINDO & SO \\
3 & execução & 1 & RUN & identificador em quatro dígitos e PID; PID acima de 9 aparece como traço \\
4 & aceite & 1 & OK & em branco \\
5 & recusa & 1 & ERR & em branco \\
6 & lista & 10 & PROG & \texttt{>} , quatro dígitos do identificador e o nome do índice \\
7 & quantidade & 10 & QUANTOS & \texttt{>01} a \texttt{>10}, iguais ao índice mais um \\
8 & arquivos & 1 & ARQUIVOS & FASE 3 \\
\bottomrule
\end{tabular}
\legend{Fonte: Autor.}
\end{table}

Os nomes da lista, na ordem do índice, são AREA, BUSCA, FAT, FIB, GCD, MEDIA, PRIMO, POT, DIG e SORT. Os quatro dígitos não saem dessa tabela: vêm do campo de identificador escrito pelo kernel. Os identificadores são 1 a 10, nos setores 20 a 29 da imagem do HD, na mesma ordem. O kernel guarda essa tabela em vetores. A leitura do setor zero ainda não preenche esses vetores.

Depois do \emph{reset}, a primeira tela é a de boas-vindas. \texttt{KEY[1]} abre o menu. Na primeira opção, a lista acompanha o índice: enquanto a confirmação ainda é zero, uma mudança de índice faz o kernel reescrever o identificador no mesmo modo, e o índice da tela é conservado. A confirmação chama \texttt{find\_program}. Se a busca falha, ou se \texttt{loadHD(36, 1536)} ou \texttt{loadROM} do setor do programa em 3072 não devolvem 0, a tela é ERR. No sucesso, o kernel mostra RUN com o identificador e PID 0 e chama \texttt{runat(3072)}. Não há \texttt{wait\_choice} antes do programa. O \texttt{HLT} armado devolve o menu. As opções 2 e 3 não carregam setor.

Na segunda opção, QUANTOS pergunta quantos programas entram na rodada. O valor é o índice mais um, de 1 a 10. O kernel então pede essa quantidade de identificadores na mesma lista. Ao terminar, mostra ERR. A preempção por quantum ainda não existe; a tela de erro é a recusa dessa opção, não uma falha de disco. Na terceira opção, a tela de arquivos avisa que criar, renomear e apagar ficam para a fase do catálogo. \texttt{KEY[1]} volta ao menu.

O menu não chama \texttt{input()} nem \texttt{output()}. As duas instruções também esperam \texttt{KEY[1]} e, enquanto esperam, a CPU não caminha pelas opções. O autoteste antigo, que fazia \texttt{output(111)}, não é chamado na entrada. Os códigos numéricos de sete segmentos que marcavam ``menu pronto'' ou ``insira o identificador'' pertenciam àquele menu por \emph{switches}. O menu atual não os emite. O sete segmentos recebe o \texttt{OUT} do programa em execução. No \texttt{tb\_menu}, a área produz base$\times$altura/2.

A reescrita do identificador na lista apaga a confirmação, porque toda escrita em 4091 o faz. Se \texttt{KEY[1]} chegasse no meio dessas poucas instruções, o aperto se perderia e a lista esperaria outro. A janela dura cerca de trinta instruções. Com o relógio da CPU em aproximadamente 5~kHz, são cerca de 6~ms, e só logo depois que uma seta mudou o índice. O tratamento de rebote do botão já espera cerca de 20~ms. Por isso o rebote não foi alongado.

A opção 1 carrega o executável, salta para a folha e volta no \texttt{HLT}. A escrita em \texttt{pc\_load} (4088) só arma quando o destino está em 3072--3839 e o HD não está congelado: grava PC+1 em \texttt{kret} e na RAM 1042 e desvia. Fora da folha a escrita é ignorada; o sentinela 99 do \texttt{tb\_lcd} continua. O \texttt{HLT} armado volta a esse PC, grava a causa 2 em 4089 e desarma. O \texttt{HLT} sem arma congela o PC e acende LEDR[16]. A leitura de 4088 devolve \texttt{epc}. \texttt{ctrl} (4090) continua sem efeito. A troca de contexto de 18 palavras, com \texttt{\$hi} e \texttt{\$lo}, o supervisor, o timer e a leitura do setor 0 não fazem parte dessa opção.

\section{Arquitetura do Kernel}

O arquivo único do kernel é dividido nas seguintes camadas lógicas:

\begin{enumerate}
    \item \textbf{shell e menu}: recebem comandos do usuário e apresentam os resultados;
    \item \textbf{serviços do kernel}: realizam gerenciamento de processos, memória, arquivos e E/S;
    \item \textbf{abstração de hardware}: concentra os acessos aos registradores mapeados em memória;
    \item \textbf{módulos do processador}: fornecem memória, disco, temporizador, interrupções e periféricos em Verilog.
\end{enumerate}

Essa separação é lógica, pois todos os serviços em C- são compilados como um único programa. Os módulos de \emph{hardware} pertencem ao projeto do processador, enquanto o código C- concentra os serviços executados pelo kernel.

O processador mantém um bit de modo controlado pelas \emph{traps} e pelo supervisor. Em modo usuário, operações privilegiadas não são executadas e as instruções \texttt{IN} e \texttt{OUT} solicitam serviços do kernel. Em modo privilegiado, kernel e supervisor acessam os mecanismos necessários à gerência do sistema. Essa distinção não constitui uma MMU e não pretende proteger o sistema contra binários manuais malformados.

\section{Gerenciamento de Processos}

O gerenciamento de processos é a função aprofundada neste projeto. O diretório contém até dezesseis programas, dez deles presentes na imagem inicial. O usuário seleciona de um a dez processos por rodada e pode escolher o mesmo programa mais de uma vez. A folha da vez e a janela de dados mantêm somente o ocupante em execução; as demais imagens ficam no disco. A opção 1 ainda não cria vários processos.

\subsection{Bloco de Controle de Processo}

Devido à ausência de estruturas em C-, cada campo do PCB é armazenado em um vetor. A posição de mesmo índice nos diferentes vetores representa um processo. A representação registrada no kernel contém:

\begin{itemize}
    \item \texttt{proc\_state}, para os códigos livre, pronto e executando;
    \item identificador numérico e setor do executável;
    \item PC lógico de retomada;
    \item área de contexto com os registradores usados pela ABI do compilador.
\end{itemize}

O PID é o índice do PCB e as bases das imagens são fixas, por isso esses campos não precisam ser armazenados. O contrato de troca preserva PC, \texttt{\$v0}, \texttt{\$t0}--\texttt{\$t9}, \texttt{\$gp}, \texttt{\$sp}, \texttt{\$fp}, \texttt{\$ra}, \texttt{\$hi} e \texttt{\$lo}. São 17 registradores e o PC, totalizando 18 palavras por processo. Esse é o conjunto mínimo suficiente para a ABI atual; novos registradores somente poderão ser usados pelo compilador após a revisão desse contrato.

O kernel cria um processo no primeiro PID livre. O contexto inicial previsto tem o PC na folha e os registradores zerados, com exceção de \texttt{\$gp=1536} e \texttt{\$sp=1791}. A opção 1 não chama \texttt{create\_process}: carrega o setor 36 em 1536 e entra em 3072 pelo prólogo do próprio programa. O setor de dados associado ao PID também é zerado na criação ou reutilização previstas. O mesmo executável pode originar mais de um processo porque cada instância possui PID, contexto e setor de dados próprios.

\subsection{Escalonador \emph{round-robin}}

A fila de prontos é representada pelo próprio vetor fixo de dez PCBs. Quando o processo atual encerra seu \emph{quantum}, o temporizador produz uma interrupção, o contexto é salvo e o processo volta ao estado pronto. O escalonador percorre circularmente as posições a partir do PID seguinte ao atual e escolhe a primeira marcada como pronta, sem manter uma fila FIFO separada.

O algoritmo abaixo é o contrato da multitarefa. A opção 1 não o executa: não há supervisor, timer nem \texttt{DISPATCH}.

\begin{enumerate}
    \item a \emph{trap} transfere o controle ao supervisor e congela o temporizador;
    \item o supervisor salva o contexto do usuário no \emph{mailbox} e grava a janela 1536--1791 no HD; o código do kernel permanece em 0--3071;
    \item o kernel copia o contexto para o PCB e marca o processo atual como pronto, caso ele não tenha finalizado;
    \item o escalonador seleciona o próximo processo pronto em ordem circular;
    \item o kernel coloca o contexto e os parâmetros do PID escolhido no \emph{mailbox} e solicita \texttt{DISPATCH};
    \item o supervisor carrega a folha a partir de 3072 e a janela do próximo PID em 1536, e restaura registradores e PC;
    \item em uma preempção, o contador de 50 instruções é reiniciado e a execução de usuário é liberada.
\end{enumerate}

O temporizador previsto contaria somente instruções concluídas de programas de usuário. Ele ainda não está no hardware. A contagem ficaria desabilitada enquanto kernel ou supervisor executam, durante operações de entrada e saída e enquanto uma transferência do HD aguarda conclusão, e só seria habilitada com mais de um processo pronto ou executando. O relógio da CPU, na síntese, usa \texttt{CPU\_DIV} 5000, cerca de 5 kHz. O \texttt{tb\_menu} encurta esse divisor só na simulação.

\subsection{Interrupções e término}

Temporizador, término, entrada, saída, cancelamento e operação privilegiada inválida são tratados como \emph{traps}. As causas são, respectivamente, 1 para fim do \emph{quantum}, 2 para \texttt{HLT}, 3 para \texttt{IN}, 4 para \texttt{OUT}, 5 para cancelamento e 6 para operação inválida; zero representa ausência de causa. Timer, \texttt{IN} e \texttt{OUT} salvam o PC da próxima instrução. \texttt{HLT}, cancelamento e operação inválida registram o PC corrente apenas para diagnóstico, pois o processo não será retomado. Se \texttt{HLT} e timer coincidirem, \texttt{HLT} possui prioridade e o kernel libera o PCB.

O supervisor sempre salva a RAM do processo antes de trazer o kernel, inclusive em \texttt{HLT} e cancelamento. Em \texttt{IN}, o hardware registra o índice do registrador de destino; em \texttt{OUT}, registra o valor-fonte. O kernel executa sua própria operação privilegiada de E/S e, no caso de entrada, copia o valor recebido para o registrador correspondente no contexto salvo. Em seguida, despacha novamente o mesmo PID, sem alterar sua posição no \emph{round-robin}; o saldo do quantum é preservado. O \texttt{HLT} não exige confirmação adicional de saída.

Em modo usuário, o contrato prevê a prioridade \texttt{KEY[0]}, \texttt{HLT}, \texttt{IN}/\texttt{OUT} e timer, nessa ordem. Esse modo privilegiado ainda não existe. No hardware atual, \texttt{IN} e \texttt{OUT} não geram \emph{trap}, e \texttt{KEY[0]} é o \emph{reset} da placa.

Na arquitetura-alvo, \texttt{KEY[0]} é um pedido de cancelamento, distinto da restauração da imagem inicial. Durante uma transferência de HD, o pedido fica pendente até o setor terminar. Em uma rodada, o kernel libera os PCBs e retorna ao menu sem apagar arquivos; dentro de um submenu, a operação incompleta é abandonada. A imagem inicial somente é restaurada ao reconfigurar ou energizar novamente a FPGA. No hardware já integrado, esse botão é o \emph{reset} ativo em nível baixo da placa. O menu em simulação não o lê, e o cancelamento por \emph{software} ainda não existe.

\section{Gerenciamento de Memória}

O processador utiliza arquitetura Harvard, com memória de instruções e memória de dados separadas. Ambas possuem 4096 palavras de 32 bits e endereçamento de 12 bits. O kernel permanece em 0--3071. A folha 3072--3839 e a janela de dados 1536--1791 são as regiões substituídas quando um programa entra.

\subsection{Memória de instruções}

A memória de instruções é gravável e possui duas portas. A porta A fornece ao processador a instrução indicada pelo PC; a porta B permite leitura e escrita pelo controlador de disco.

Os endereços de 0 a 3071 formam o kernel residente. Após o \emph{reset}, \texttt{so.txt} traz o kernel compilado nas palavras 0 a 1334, NOP de 1335 a 3070 e o salto de guarda \texttt{J 3071} na palavra 3071. A folha da vez, 3072 a 3839, recebe um programa. Os \texttt{J} e \texttt{JAL} do alvo \texttt{os-user} já somam 3072, e \texttt{BEQ} e \texttt{BNE} continuam relativos. Os endereços de 3840 a 4095 estão reservados ao supervisor protegido, que ainda não existe e permanece em zero. O kernel não é descartado para o programa executar. A porta B continua recebendo setores do disco.

\subsection{Memória de dados}

Os endereços de 0 a 1023 são dados e pilha do kernel, quatro setores. O \emph{reset} do hardware põe \texttt{\$sp} em 1023; o alvo padrão do compilador emite \texttt{addi \$sp,\$zero,1015}. Um programa de usuário usa a janela 1536 a 1791, com \texttt{\$gp = 1536} e \texttt{\$sp = 1791}. A região de 1024 a 1279 guarda o estado protegido e o \emph{mailbox}. O contrato de 18 palavras, PC e 17 registradores incluindo \texttt{\$hi} e \texttt{\$lo}, continua especificado e não é o que a opção 1 grava. \texttt{runat} escreve \texttt{\$v0}, \texttt{\$t0} a \texttt{\$t9}, \texttt{\$gp}, \texttt{\$sp}, \texttt{\$fp} e \texttt{\$ra} em 1043 a 1057; \texttt{\$hi} e \texttt{\$lo} ficam de fora, e o hardware grava PC+1 em 1042. A região de 1280 a 1535 oferece um buffer de setor; 1792 a 4087 permanece reservada; e 4088 a 4095 pertence à E/S mapeada em memória.

O contrato inicial do \emph{mailbox} reserva 1024--1041 para o contexto de usuário, 1042--1059 para o contexto do kernel, 1060 para a causa, 1061 para o PID atual, 1062 para o próximo PID, 1063 para o destino de \texttt{IN}, 1064 para o valor de \texttt{OUT}, 1065 e 1066 para os setores de código e dados do próximo PID, 1067 como reservado e 1068 para o saldo do \emph{quantum}. Esse despacho por \texttt{ctrl} ainda não tem efeito no hardware, e o modo privilegiado não está implementado. A opção 1 usa o \texttt{SW} em 4088, não o supervisor. \texttt{runat} grava quinze registradores em 1043--1057 e deixa \texttt{\$hi} e \texttt{\$lo} de fora.

O registrador \texttt{\$gp} do kernel permanece em zero. No alvo \texttt{os-user} o prólogo é \texttt{addi \$gp,\$zero,1536} e \texttt{addi \$sp,\$zero,1791}, depois \texttt{j main}. O alvo padrão continua com \texttt{addi \$sp,\$zero,1015} e saltos a partir de 0. Não há realocação de dados dentro da janela: o programa já nasce com esses endereços.

\subsection{Política de transferência}

Antes de trocar o ocupante, o contrato prevê que uma \emph{trap} ou chamada de despacho entre no supervisor. Ele salvaria o contexto, gravaria no HD a janela do ocupante, substituiria só a folha a partir de 3072 e restauraria PC e registradores. O controlador transfere setores de 256 palavras, uma palavra por ciclo de \texttt{clk\_cpu}. A opção 1 já faz duas transferências sem supervisor: o setor 36 para 1536 e o setor do programa para 3072. Não recarrega o kernel.

Código de kernel e executáveis possui cópia prevista no HD e não é automodificável, portanto não há escrita de retorno da IMEM. O retorno da opção 1 não recarrega setores de código: o \texttt{HLT} armado desvia para o PC salvo no kernel, que permaneceu em 0--3071. A troca completa da janela e o supervisor continuam pendentes. A RAM mutável, no contrato completo, continua sendo gravada antes de ser substituída.

A primeira versão não acrescenta MMU, verificação de limites ou tratamento de instruções privilegiadas malformadas. O objetivo é validar os algoritmos de gerenciamento com os programas produzidos pelo compilador. Essas proteções permanecem como possíveis extensões, não como requisitos da implementação avaliada.

\subsection{Contratos do compilador}

O compilador produz o alvo padrão e o alvo \texttt{os-user}. O alvo \texttt{os-kernel} ainda não existe na linha de comando. O alvo padrão emite \texttt{addi \$sp,\$zero,1015} e saltos a partir de 0. Ele gera a ROM solta de \texttt{exemplos/} e também o kernel gravado em \texttt{so.txt}. O alvo \texttt{os-user} produz o prólogo \texttt{addi \$gp,\$zero,1536} e \texttt{addi \$sp,\$zero,1791}, depois \texttt{j main}. \texttt{J} e \texttt{JAL} somam 3072. \texttt{BEQ} e \texttt{BNE} continuam relativos. Imagem acima de 256 palavras é recusada, assim como chamadas de HD, LCD e \texttt{runat}. Os programas do menu saem desse alvo, nos setores 20 a 29.

Os executáveis de usuário e de teste começam com um NOP cujo campo de endereço informa o tamanho total do código e recebem \texttt{HLT} automático. Um programa de usuário ocupa no máximo um setor, incluindo o NOP inicial e o \texttt{HLT} final. O kernel compilado em \texttt{so.txt} tem 1335 palavras, também começa com o NOP de tamanho e termina o corpo com \texttt{HLT}. O \texttt{J 3071} da palavra 3071 é o guarda da imagem, fora dessas 1335 palavras. A função \texttt{main} do menu não retorna. O metadado de tamanho não reduz a transferência fixa de 256 palavras.

O alvo \texttt{os-user} já recusa código acima de 256 palavras. O alvo \texttt{os-kernel}, ainda ausente, deverá rejeitar código que não caiba antes do salto de guarda na palavra 3071 e dados estáticos incompatíveis com a RAM 0--1023. A profundidade dinâmica da pilha será verificada nos testes de integração.

Para que o kernel acesse os registradores mapeados e o buffer do HD sem aritmética geral de ponteiros, o alvo \texttt{os-kernel} oferecerá os intrínsecos \texttt{mmio\_read(addr)}, \texttt{mmio\_write(addr, value)}, \texttt{mem\_read\_abs(addr)} e \texttt{mem\_write\_abs(addr, value)}. O endereço deverá ser um literal constante: \texttt{mmio\_*} aceitará somente 4088--4095 e \texttt{mem\_*}, 0--4095. Eles serão traduzidos para \texttt{LW}/\texttt{SW} com base \texttt{\$zero} e endereço imediato absoluto de 14 bits e serão rejeitados em \texttt{os-user}. Esse alvo ainda não existe.

\texttt{void runat(int)} já existe no alvo padrão e não é \texttt{jal}. Grava \texttt{\$v0}, \texttt{\$t0}--\texttt{\$t9}, \texttt{\$gp}, \texttt{\$sp}, \texttt{\$fp} e \texttt{\$ra} em 1043--1057, depois o \texttt{SW} do endereço em 4088, depois os \texttt{LW} na mesma ordem, começando em \texttt{lw \$v0,1043}. \texttt{\$hi} e \texttt{\$lo} ficam de fora. O alvo \texttt{os-user} recusa essa chamada.

\section{Sistema de Arquivos}

O sistema de arquivos armazena programas executáveis em até oito diretórios planos, sem subdiretórios. O disco simulado possui 8 trilhas, 8 setores por trilha e 256 palavras de 32 bits por setor. Sua capacidade total é de 64 setores, 16.384 palavras ou 64 KiB.

O software endereça o disco por LBA entre 0 e 63. Internamente, a tradução utiliza \texttt{track = sector\_lba / 8}, \texttt{sector\_in\_track = sector\_lba \% 8} e \texttt{flat\_addr = sector\_lba * 256 + word\_offset}.

O setor zero está especificado como catálogo durável. As seis primeiras palavras guardariam assinatura, versão, geração, capacidade de diretórios, capacidade de arquivos e soma de verificação. Esta última é o complemento de dois da soma das demais 255 palavras, de modo que a soma do setor inteiro resulte em zero. Em seguida existem oito entradas de diretório e dezesseis entradas de arquivo. O setor do executável é derivado da posição de sua entrada, entre os setores 20 e 35. O setor 46 está especificado como cópia-espelho. A opção 1 não lê esse catálogo: usa vetores fixos e os programas dos setores 20 a 29.

Como a linguagem C- não possui cadeias de caracteres, nomes de diretórios e IDs de arquivos são números entre 1 e 1023. Eles pertencem a espaços de nomes separados, de modo que um diretório e um arquivo podem usar o mesmo número. IDs de arquivo são globalmente únicos. Arquivos guardam o slot interno do diretório, e não seu nome; assim, renomear um diretório não exige alterar suas entradas de arquivo. O valor zero representa o SO ou a ação de cancelar e voltar.

\begin{table}[H]
\centering
\caption{Organização do disco simulado}
\label{tab:layout-disco}
\small
\begin{tabular}{p{3.2cm}p{9.2cm}}
\toprule
\textbf{Setores} & \textbf{Conteúdo} \\
\midrule
0 & Catálogo persistente de diretórios e programas \\
1--12 & Código do kernel, 3072 palavras, com o salto de guarda na palavra 3071. Na imagem atual esses setores estão em NOP; o kernel compilado está em \texttt{so.txt} \\
13--16 & Dados e pilha do kernel \\
17--19 & Reservados \\
20--35 & Até 16 programas executáveis, um por setor \\
36--45 & Imagens de RAM dos PIDs 0 a 9 \\
46 & Cópia-espelho transacional do catálogo \\
47--63 & Até 17 pacotes de executáveis preparados para importação \\
\bottomrule
\end{tabular}
\legend{Fonte: Autor.}
\end{table}

No estado inicial, somente o diretório raiz está ativo. Ele ocupa o slot zero, começa com o nome 1 e contém os dez programas iniciais. A raiz pode ser renomeada, mas não excluída. Os demais sete diretórios podem ser criados, selecionados, renomeados e excluídos. O kernel mantém um diretório atual, e as operações de arquivo atuam somente sobre ele.

As operações de diretório são listar, criar, selecionar, renomear e excluir. Excluir um diretório exige confirmação explícita e desativa logicamente todos os arquivos nele contidos; os setores deixam de estar referenciados e podem ser reutilizados. As operações de arquivo já definidas são:

\begin{itemize}
    \item \textbf{listar}: percorrer as entradas ativas do diretório atual;
    \item \textbf{carregar}: localizar o programa no diretório atual e transferir seu código para a IMEM;
    \item \textbf{importar}: validar um pacote preparado, copiá-lo para o setor de um slot livre e ativar sua entrada no catálogo;
    \item \textbf{renomear}: substituir seu ID por outro ID globalmente livre;
    \item \textbf{mover}: alterar somente o slot de diretório armazenado na entrada;
    \item \textbf{remover}: desativar a entrada, sem compactar o catálogo.
\end{itemize}

O conteúdo de um executável não é editado depois que ele entra no sistema de arquivos. A importação utiliza os pacotes posicionados pelo empacotador nos setores 47 a 63. O usuário seleciona o número do pacote, o diretório de destino e um ID globalmente livre. O kernel carrega o setor escolhido no buffer, valida o tamanho registrado no NOP inicial e a presença do \texttt{HLT} final e calcula uma soma de verificação das 256 palavras. Em seguida, escolhe a primeira entrada de arquivo livre, grava o setor completo no setor 20 mais o índice dessa entrada, relê o destino e compara sua soma com a origem. Somente depois prepara uma nova geração do catálogo que torna o arquivo visível. Pacotes vazios, inválidos ou maiores que 256 palavras são rejeitados sem alterar o catálogo.

No contrato do catálogo, ainda não executado, o kernel carregaria os setores zero e 46, validaria assinatura, versão, capacidades e soma de verificação e escolheria a cópia válida de maior geração. Se apenas o espelho for válido ou estiver mais atualizado, ele é copiado de volta ao setor zero antes da abertura do menu. Em cada alteração, o kernel monta a versão candidata no buffer sem modificar os vetores ativos, incrementa sua geração, calcula a soma de verificação e grava primeiro o setor 46. A releitura e validação do espelho constituem o ponto de \emph{commit}: a partir dele, a mudança é publicada nos vetores. O kernel tenta então sincronizar e verificar o setor zero. Se essa segunda escrita falhar, a operação permanece confirmada pelo espelho, o usuário recebe um aviso de recuperação pendente e uma nova alteração de catálogo é recusada até que a sincronização seja concluída. Um cancelamento anterior ao \emph{commit} descarta o candidato; depois dele, não desfaz a mudança. Se houver perda de energia entre as duas cópias, a inicialização recupera a geração válida mais recente.

Cada processo usa o setor de dados correspondente ao PID, por isso duas instâncias do mesmo executável preservam dados e contexto independentes. Alterações cujo \emph{commit} terminou persistem durante cancelamentos controlados; uma nova energização ou reconfiguração da FPGA restaura a imagem inicial.

\section{Gerenciamento de Entrada e Saída}

O kernel organiza logicamente a entrada e saída em serviços e controladores dentro do arquivo único. O menu utiliza entrada numérica confirmada por botão, saída numérica e LCD; o disco utiliza registradores próprios mapeados em memória.

O disco e os controles do sistema usam E/S mapeada em memória. Os endereços reservados e suas funções são apresentados na \autoref{tab:mmio}.

\begin{table}[H]
\centering
\caption{Registradores de E/S mapeados em memória}
\label{tab:mmio}
\small
\begin{tabular}{clp{7.8cm}}
\toprule
\textbf{Endereço} & \textbf{Nome} & \textbf{Função} \\
\midrule
4088 & \texttt{pc\_load} & Escrita com destino em 3072--3839 e HD descongelado grava PC+1 em \texttt{kret} e na RAM 1042 e desvia. Fora da folha a escrita é ignorada. A leitura devolve \texttt{epc}. \\
4089 & \texttt{cause} & O \texttt{HLT} armado grava 2 e desarma. Timer, \texttt{IN}, \texttt{OUT} e cancelar continuam especificados e ainda não são produzidos. \\
4090 & \texttt{ctrl} & Especificado para modo, temporizador e \texttt{DISPATCH}. No hardware integrado a escrita não tem efeito. \\
4091 & \texttt{display\_ctx} & Modo, identificador e PID na escrita. A leitura acrescenta o índice da opção e o bit 27 de confirmação. O catálogo está na \autoref{tab:lcd}. \\
4092 & \texttt{hd\_cmd} & 0: nenhuma operação; 1: HD$\rightarrow$RAM; 2: HD$\rightarrow$IMEM; 3: RAM$\rightarrow$HD. Não há comando de IMEM para o HD. \\
4093 & \texttt{hd\_sector} & Seleciona o setor lógico. \\
4094 & \texttt{hd\_addr} & Define o endereço-base na memória. \\
4095 & \texttt{hd\_status} & Somente leitura: 0 ocioso, 1 ocupado, 2 concluído e 3 erro. \\
\bottomrule
\end{tabular}
\legend{Fonte: Autor.}
\end{table}

Para iniciar uma transferência, o kernel escreve o setor em 4093, o endereço-base em 4094 e o comando em 4092. O controlador captura esses três parâmetros uma única vez em registradores internos, ignora novas escritas de início enquanto ocupado, congela PC e sinais comuns de escrita e transfere uma palavra por ciclo. Após aposentar a 256ª palavra, encerra o comando, sinaliza conclusão, libera a CPU e permite que o PC avance exatamente uma vez. O intrínseco, já com o PC liberado, lê \texttt{hd\_status} e devolve 0 quando o estado é concluído e 1 quando é erro.

Na leitura de \texttt{ctrl}, o contrato reserva os bits 0 a 4 para modo privilegiado, temporizador habilitado, \emph{trap} pendente, HD ocupado e cancelamento pendente. Em uma escrita privilegiada, os bits 8, 9, 10 e 11 seriam pulsos de comando para habilitar o temporizador, desabilitá-lo, executar \texttt{DISPATCH} e reconhecer a causa tratada. Esse mapa ainda não atua: o endereço 4090 não altera o processador.

O estado do HD passa de 0 para 1 quando um comando válido é aceito e termina em 2 ou 3. O valor 2 indica que as 256 palavras foram transferidas; o valor 3 informa setor, endereço ou comando inválido. Ao retomar, o intrínseco devolve 0 em sucesso e 1 em erro, e a opção 1 trata o valor diferente de zero como \texttt{show\_error}. Alterações de catálogo não são publicadas antes do ponto de \emph{commit} definido no sistema de arquivos, contrato que a opção 1 ainda não exerce: o setor 0 não é lido.

O compilador converte \texttt{input()} e \texttt{output()} nas instruções \texttt{IN} e \texttt{OUT}. No processador integrado elas acessam o periférico direto e não geram \emph{trap}. O \texttt{tb\_menu} confere que o \texttt{KEY[1]} de um \texttt{IN} não confirma o menu. O contrato ainda pendente é que, em modo usuário, essas instruções registrem causa, PC de retomada e operando e transfiram o controle ao supervisor. Esse modo privilegiado não está implementado.

No contrato ainda pendente, ambas as operações aguardam confirmação em \texttt{KEY[1]}, bloqueiam o sistema inteiro e congelam o temporizador. Hoje o \texttt{IN} e o \texttt{OUT} já esperam \texttt{KEY[1]} no próprio programa. Cada \texttt{OUT} pausa uma única vez; o \texttt{HLT} seguinte não solicita outra confirmação e o último valor permanece nos sete segmentos. O LCD utiliza as mensagens fixas da \autoref{sec:menu-lcd}. Durante um programa, a tela RUN mostra o identificador e o PID. A caminhada do menu não usa \texttt{IN} nem \texttt{OUT}.

O bloqueio global durante \texttt{IN} e \texttt{OUT}, com o supervisor retirando a imagem do processo, permanece especificado e não está implementado. A opção 1 deixa o programa na folha enquanto o \texttt{IN} espera \texttt{KEY[1]}. Escalonar outro usuário enquanto se aguarda uma ação humana exigiria transformar a chamada em uma requisição pendente. Uma versão posterior poderá introduzir o estado bloqueado no PCB.

\section{Construção das Imagens Iniciais}

A imagem de \emph{reset} já usada pelo Quartus é \texttt{processador/so.txt}, lida com \texttt{hd\_image.txt} pelo \texttt{MISC\_FILE} do \texttt{.qsf}. \texttt{so.txt} tem 4096 palavras: o kernel nas palavras 0 a 1334, NOP de 1335 a 3070, \texttt{J 3071} na palavra 3071 e zero de 3072 a 4095. O supervisor não está nessa imagem. Os padrões de síntese são \texttt{CPU\_DIV} 5000, \texttt{BTN\_STABLE} 1000000 e os tempos reais do LCD. O \texttt{tb\_menu} encurta esses parâmetros só na simulação. Os dez programas canônicos ocupam os setores 20 a 29 de \texttt{hd\_image.txt}. \texttt{compilador/exemplos/} não é essa imagem: \texttt{reset\_tests.sh} sem flag regenera a ROM solta e, com \texttt{--target=os-user}, grava só os setores 20 a 29.

Um empacotador determinístico ainda deverá fechar a imagem final, inclusive o supervisor em 3840--4095, que não será compilado pelo alvo C-. O salto de guarda já ocupado é o endereço 3071, com \texttt{J 3071}. Os setores 46 a 63 e a leitura do setor 0 continuam fora da opção 1.

\section{Inicialização e Fluxo de Execução}

O fluxo previsto de inicialização é o seguinte. Só o primeiro e o quarto itens ocorrem hoje. Os itens do catálogo continuam por fazer, porque o setor 0 não é lido:

\begin{enumerate}
    \item ao ligar a placa, \texttt{so.txt} coloca o kernel no endereço zero; o supervisor previsto permanece ausente;
    \item os setores zero e 46 seriam validados e a geração mais recente do catálogo seria recuperada nos vetores;
    \item a tabela de processos seria inicializada;
    \item o menu é apresentado e aguarda uma opção do usuário.
\end{enumerate}

Na execução individual já simulada, a opção 1 localiza o programa, carrega o setor 36 em 1536 e o executável em 3072, mostra RUN e chama \texttt{runat(3072)}. O \texttt{HLT} armado volta ao kernel, que não foi retirado da IMEM. Não há supervisor nem desvio para o endereço zero.

No contrato de multitarefa, ainda não ligado ao menu, os programas selecionados entrariam na fila de prontos. A cada interrupção o supervisor retornaria ao kernel, que atualizaria o PCB e selecionaria o próximo processo pronto. Esse fluxo não está implementado.

Se \texttt{KEY[0]} for acionado enquanto um programa ocupa a folha, o contrato prevê que o cancelamento aguarde uma eventual transferência de setor terminar. No hardware já integrado, esse botão é o \emph{reset} da placa. O supervisor que recarregaria o kernel ainda não existe. O catálogo e os executáveis do HD não são apagados por esse \emph{reset} de software, que também ainda não existe.

\section{Técnicas e Algoritmos Selecionados}

A \autoref{tab:decisoes} resume as principais decisões do sistema.

\begin{table}[H]
\centering
\caption{Técnicas e algoritmos definidos para o projeto}
\label{tab:decisoes}
\small
\begin{tabular}{p{3.0cm}p{4.0cm}p{5.7cm}}
\toprule
\textbf{Área} & \textbf{Método} & \textbf{Implementação} \\
\midrule
Kernel & Monolítico; supervisor ainda ausente & Kernel em um arquivo C-. O Assembly de 3840--4095 não está na imagem. \\
Processos & \emph{Round-robin} previsto & A opção 1 não cria processo nem preempta. Timer e busca circular continuam pendentes. \\
Contexto & ABI de 18 palavras prevista & A opção 1 grava 15 registradores em 1043--1057, sem \texttt{\$hi} e \texttt{\$lo}. \\
Memória & Folha e janela substituíveis & O kernel fica em 0--3071. A opção 1 troca a folha 3072--3839 e carrega o setor 36 em 1536. \\
Arquivos & Oito diretórios planos e setores fixos & Dezesseis executáveis previstos; o catálogo fixo da opção 1 não lê o setor 0. \\
Disco & MMIO e transferência bloqueante & Comandos 1, 2 e 3 em 4092--4095; 256 palavras por setor; retorno 0 ou 1. \\
E/S & \texttt{IN}/\texttt{OUT} diretos nesta etapa & Esperam \texttt{KEY[1]} no programa. A \emph{trap} que traria o kernel continua pendente. \\
Inicialização & Kernel presente no \emph{reset} & \texttt{so.txt} já contém o kernel. A volta da opção 1 é o \texttt{HLT} armado, não o supervisor. \\
\bottomrule
\end{tabular}
\legend{Fonte: Autor.}
\end{table}

\chapter{Conclusão}

O estudo relaciona as funções de gerenciamento de processos, memória, arquivos e entrada e saída às restrições do processador MIPS-Lite e da linguagem C-. A ausência de ligação entre múltiplos arquivos conduziu à especificação de um kernel monolítico em \texttt{so.cms}. O kernel permanece residente em 0--3071 enquanto a folha recebe o programa. O supervisor Assembly em 3840--4095 continua especificado e ausente: não constitui BIOS, processo ou segundo kernel, e não devolve o controle na opção 1.

O desenho ainda prevê \emph{round-robin}, o contexto de 18 palavras com \texttt{\$hi} e \texttt{\$lo}, e a gravação da janela de dados. A opção 1 já executa um programa sem esses mecanismos. IMEM e RAM têm 4096 palavras. O disco tem 64 setores de 256 palavras. Os dez executáveis canônicos estão nos setores 20 a 29. Os setores 1 a 12 do HD seguem em NOP. O setor 0 ainda não é lido.

O sistema de arquivos foi definido com oito diretórios planos e operações para listar, importar, carregar, renomear, mover e excluir programas. Essas operações não estão ligadas ao menu. A E/S dos programas por \emph{trap} permanece especificada. No hardware atual, \texttt{IN} e \texttt{OUT} acessam o periférico direto.

O menu do LCD é percorrido em simulação, e a opção 1 carrega o programa. O \texttt{tb\_menu} executa a área, com dois \texttt{IN} e o \texttt{OUT} de base$\times$altura/2, volta ao menu com \texttt{\$gp} 0, carrega um segundo setor na folha e confere que o \texttt{KEY[1]} de um \texttt{IN} não confirma o menu. \texttt{loadHD}, \texttt{storeHD} e \texttt{loadROM} já usam 4092 a 4095. Continuam pendentes o temporizador, o modo privilegiado, o supervisor, a troca de 18 palavras, o \emph{round-robin}, o CRUD, a leitura do setor 0 e \texttt{IN}/\texttt{OUT} como \emph{trap}.

As principais dificuldades previstas são respeitar os limites de código e pilha do C-, manter a ABI sincronizada com o contexto salvo, garantir que cada transferência bloqueante avance o PC uma única vez e diagnosticar falhas que atravessam compilador, software e Verilog. O custo do \emph{swapping} e o bloqueio global durante \texttt{IN}/\texttt{OUT} também limitam o desempenho, embora simplifiquem a primeira implementação. Essas limitações deverão ser medidas nos testes e poderão orientar versões posteriores com processos bloqueados, E/S concorrente e proteção de memória.

Essas definições estabelecem a organização da arquitetura-alvo e sua comunicação com o compilador e o processador, mas sua conformidade dependerá dos resultados experimentais das próximas etapas.

% ELEMENTOS PÓS-TEXTUAIS
\phantompart
\postextual

\bibliography{referencias}

\phantompart
\printindex

\end{document}
